\documentclass{article}
\usepackage{amsmath}
\usepackage{amsthm}

\newtheorem{theorem}{Problem}

\begin{document}

\begin{theorem}
What is the sum of all of the multiples of 7 between 100 and 200?
\end{theorem}

\begin{proof}
Step 0: I think that we can solve this problem by using the arithmetic series formula.
Step 1: That sounds like a good idea. I remember that the arithmetic series formula is $S_n = \frac{n}{2}(2a_1 + (n-1)d)$ where $n$ is the number of terms, $a_1$ is the first term, and $d$ is the common difference.
Step 2: I remember that too. Now let's use that formula to solve the problem.
Step 3: Ok. First we need to find $n$, the number of terms.
Step 4: That's right. So let's start by listing all of the multiples of 7 between 100 and 200.
Step 5: They are 105, 112, 119, 126, 133, 140, 147, 154, 161, 168, 175, 182, 189, 196.
Step 6: Now let's count them.
Step 7: There are 14 multiples of 7 between 100 and 200.
Step 8: So $n = 14$. We also need to find the first term. That would be 105.
Step 9: And the common difference is 7, since we are adding 7 each time to get the next term.
Step 10: So let's plug those numbers into the formula. $S_n = \frac{n}{2}(2a_1 + (n-1)d)$ becomes $S_{14} = \frac{14}{2}(2(105) + (14-1)*7)$.
Step 11: Now let's simplify. $\frac{14}{2}$ is 7, and $(14-1)*7$ is 91, so we now have $S_{14} = 7(210 + 91)$.
Step 12: Yes, and 210 + 91 is 300. So $S_{14} = 7 * 300$.
Step 13: And $7 * 300$ is 2100.
Step 14: So the sum of all the multiples of 7 between 100 and 200 is 2100.
Step 15: That's the final answer.
\end{proof}

\end{document}
